% Section 8 (part 1 of 8): the connection problem. Label prefix: con.
% No new macros. Letters used in one proof only are declared there and are local to it.

\section{The connection problem}
\label{sec:con}

The canonical gauge $\Pc$ of \cref{sec:gau} is normalised at the conifold point $u=0$, the rows of the formal quintic at the large-radius point $q=0$. The normalisation series $\Dn_0$ of \cref{thm:gau-R0} measures the difference, and $\Fg{g}$ for every $g$ (\cref{def:gau-Fg}) is the statement $\Dn_0=1$. This section reduces $\Dn_0=1$ to one identity at $q=0$. After the polynomial structure of the colour-$0$ WKB functions (\cref{sec:con-problem}) we define the full-Stirling saddle series $c(z;L)$ of a Mellin--Barnes summand and the hatted series $\PhiSP$ (\cref{sec:con-form}). Theorem~P (\cref{sec:con-P}) identifies $\PhiSP$ with the canonical function $\Phican$: part~(a) shows that $\PhiSP$ solves HPF and has polynomial coefficients, part~(b) that it is the canonical one, by an exact identity at the conifold end, the residue of the $1$-form at its pole. The glue (\cref{sec:con-glue}) evaluates the comparison \cref{thm:gau-R0} at $u=1$ and gives $\Dn_0(z)=\PhiSP(z,1)/h(z)$ exactly; the value $\PhiSP(z,1)=c(z;1)$ is determined in \cref{sec:gam}. We work in colour~$0$; the other colours follow by \cref{cor:gau-colours}. The formal stationary-phase lemmas are those of \cref{sec:fsp}. Topological recursion is not used in this section. No quoted input enters \cref{thm:con-Pa}, \cref{prop:con-deg} or \cref{thm:con-Pb}; the last concerns the cyclic system of a conifold frame and its canonical gauge, not the rows of the formal quintic. The glue \cref{thm:con-glue} uses the input (QRR), once, and the actual rows, which rest on (M) and (G0).

\subsection{The normalisation problem and the polynomial structure}
\label{sec:con-problem}

\emph{Conventions.} A colour-$0$ WKB function is a series $\Phi(\zh,u)=\sum_{n\geq 0}\mathrm{ph}_n(u)\zh^n$ with $\mathrm{ph}_0=1$; its flat amplitude is $A_n:=L^{4n+1}\mathrm{ph}_n(L^{-5})$, so that $\sum_n A_nz^n=L\,\Phi(L^4z,L^{-5})$ (\cref{def:set-points,def:set-rows}). We write $\PFfive^{L}$ for $\PFfive$ with $z\Dq$ replaced by $z\Dq+L$, that is, $\PFfive^{L}=\re^{-u_0/z}\,\PFfive\,\re^{u_0/z}$ with $\Dq u_0=L$; it is the operator $\PFfive^{(0)}$ of \cref{lem:gau-HPF,lem:gau-wkb}. At fixed $z$ one has $u\,\dd/\dd u=\theta_u-\tfrac45\theta_{\zh}$ on functions of $(\zh,u)$ and $z/u=\zh/s$, so that for $e\in\bbZ$ and $F$ in $\bbQ\lser{s}\fps{\zh}$
\begin{equation}
\label{eq:con-chain}
(z\Dq+L)(s^eF)=s^{e-1}\mathscr O_eF,\qquad \mathscr O_e=1+\zh(1-u)\bigl(\tfrac45\theta_{\zh}-\tfrac e5-\theta_u\bigr),
\end{equation}
the chain rule of \cref{lem:gau-HPF}. Applying it five times to $s^{-1}\Phi$, with $\PFfive$ written as
\[
\PFfive=u(z\Dq)^5-1-(1-u)\sum_{m=1}^5e_m\bigl(\tfrac15,\dots,\tfrac55\bigr)\,z^m(z\Dq)^{5-m}
\]
($e_m$ the elementary symmetric functions), gives $\PFfive^{L}(s^{-1}\Phi)=s^{-1}\,\mathrm{HPF}(\Phi)$, where $\mathrm{HPF}(\Phi)$ is the left side of \eqref{eq:gau-HPF}; so $L\re^{u_0/z}\Phi$ solves $\PFfive$ if and only if $\Phi$ solves HPF, as identities in $\bbQ\lser{s}\fps{\zh}$. The free coefficients of $\Phi$ are $\fc_m=[u^{4m}]\mathrm{ph}_{5m}$ (\cref{lem:gau-HPF}).

\begin{rmk}[The normalisation problem]
\label{rmk:con-problem}
The canonical gauge is normalised at the conifold point, by the gauge identity (G) (\cref{thm:gau-G,def:gau-can}); the rows of the formal quintic are normalised at $q=0$, by (QRR) (\cref{thm:set-qrr}). The series $\Dn=\diag(\Dn_i)$ of \cref{thm:gau-R0,cor:gau-colours} compares the two, and Property N($g$) (\cref{def:set-N}) needs $\Fg{g}$ (\cref{def:gau-Fg}), through the comparison theorem \cref{thm:flh-main} and \cref{thm:mai-main}. The series $\Dn$ is a connection constant: it compares a series built at $u=0$ with its value at $u=1$. It is not determined by data at $q=0$ alone, nor by data at $u=0$ alone (\cref{lem:con-u0,rmk:con-model}).
\end{rmk}

\begin{rmk}[Genus bookkeeping]
\label{rmk:con-genus}
By \cref{def:gau-Fg}, $\Fg{g}$ is the vanishing of the coefficients of $z,\dots,z^{3g-3}$ in $\Dn_0-1$; and $\Dn_0\equiv 1\bmod z^{N+1}$ if and only if these coefficients vanish through $z^N$; the coefficient of $z^n$ enters $\Fg{g}$ first at $g=\lceil(n+3)/3\rceil$. The identity $\Dn_0(z)\Dn_0(-z)=1$ would halve the conditions: $\log\Dn_0=\sum_na_nz^n$ would be an odd series, and $\Dn_0=1$ would follow from $a_n=0$ for every odd $n$. This identity is equivalent to the symplecticity of the rows expanded in the frame. Indeed, the model columns satisfy $\Qr^{\mathrm m}(-z)^T\eta^{-1}\Qr^{\mathrm m}(z)=1$, because $\Mo(-t)^T\Mo(t)=1$, $\Psi_0^T\eta^{-1}\Psi_0=1$ (\cref{def:set-M0,lem:set-M0sympl}) and the powers of $\varsigma$ cancel on the pairs $(a,3-a)$, where $\nu_a+\nu_{3-a}=2$; and $\Pc(-z)^T\eta\Pc(z)=\eta$ (\cref{prop:gau-symp}), with $\eta^{-1}=\eta/25$. Inserting $\Qr=\Pc\Qr^{\mathrm m}\Dn^{-1}$ (\cref{cor:gau-colours}) gives $\Qr(-z)^T\eta^{-1}\Qr(z)=(\Dn(-z)\Dn(z))^{-1}$, colour by colour; equivalently, the outer matrix $R'$ is symplectic if and only if $\Dn_0(z)\Dn_0(-z)=1$ (\cref{lem:flh-inputs}~(iii)). The symplecticity of $R$ (\cref{def:set-rows}~(a)) is an identity between power series in $q$; to carry it to the expanded rows one would need it as an identity between the representatives of the rows in $\Nring\otimes\bbQ(\zeta)$ (\cref{cor:gau-rowsNN}), for instance from the injectivity of $\ev$ on them, which is not claimed (see the paragraph after \cref{lem:flh-inputs}). So the halving is not available a priori, and it is not used: $\Dn_0=1$ is proved directly (\cref{thm:gam-D}), and then $\Dn_0(z)\Dn_0(-z)=1$ holds trivially. The finite ranges in which $\Dn_0=1$ has been tested are reported in \cref{sec:evi-normal}; they are not used.
\end{rmk}

\begin{lem}[The conifold value]
\label{lem:con-u0}
Let $\Phi=\sum_n\mathrm{ph}_n(u)\zh^n$ solve HPF, with $\mathrm{ph}_0=1$ and every $\mathrm{ph}_n$ in $\bbQ\fps{u}$. Then $\mathrm{ph}_n(0)=[t^n]\Fs_0$ for every $n\geq 1$, whatever the free coefficients; that is, $\Phi(\zh,0)=\Fs_0(\zh)$.
\end{lem}

\begin{proof}
Put $\Phi_0(\zh):=\Phi(\zh,0)$ and $\mathscr O^0_e:=1+\zh(\tfrac45\theta_{\zh}-\tfrac e5)$, the operator $\mathscr O_e$ at $u=0$. The operators $\theta_u$ and multiplication by $u$ do not lower the $u$-order and $\theta_u$ kills $u^0$; the terms of HPF in the bracket of \eqref{eq:gau-HPF} other than $3\zh\,\Phi^{[4]}$ carry a factor $u^k$, $k\geq 1$. So the $u^0$-part of HPF involves only $\Phi_0$ and reads
\[
(\mathscr O^0_{-5}-3\zh)\,\mathscr O^0_{-4}\mathscr O^0_{-3}\mathscr O^0_{-2}\mathscr O^0_{-1}\Phi_0=\Phi_0 .
\]
By \cref{lem:gau-HPF} its $\zh^{n+1}$-coefficient is $4n\,\mathrm{ph}_n(0)+r_n(0)=0$, where $r_n(0)\coloneqq\coef{u^0}r_n$ involves only the values at $u=0$ of $\mathrm{ph}_0,\dots,\mathrm{ph}_{n-1}$; the pivot $4n$ is non-zero for $n\geq 1$, so this equation has exactly one solution with constant term $1$. This equation is the model's. Put $g_e:=(4\theta_{\zh}-e)/5$; since $\theta_{\zh}\zh=\zh(\theta_{\zh}+1)$, one has $g_e\zh=\zh g_{e-4}$ and
\[
\mathscr O^0_a\mathscr O^0_b\mathscr O^0_c=1+\zh\,(g_a+g_b+g_c)+\zh^2\,(g_{a-4}g_b+g_{a-4}g_c+g_{b-4}g_c)+\zh^3\,g_{a-8}g_{b-4}g_c .
\]
For $(a,b,c)=(10,-4,-3)$ and for $(a,b,c)=(0,1,2)$ both sides equal $1+\zh(12\theta-3)/5+\zh^2(48\theta^2+24\theta-18)/25+\zh^3(4\theta-2)(4\theta+8)(4\theta+3)/125$, $\theta:=\theta_{\zh}$. Since $\mathscr O^0_{-5}-3\zh=\mathscr O^0_{10}$, the $u^0$-part of HPF is
\[
\Phi_0=\mathscr O^0_0\mathscr O^0_1\mathscr O^0_2\mathscr O^0_{-2}\mathscr O^0_{-1}\Phi_0 .
\]
In colour $0$ the operators of the T-system of \cref{prop:set-modrows} are $\mathcal T_k=-\mathscr O^0_k$ at $t=\zh$, so the closing relation of that system, $\Fs_0=-\mathcal T_0\mathcal T_1\mathcal T_2\mathcal T_{-2}\mathcal T_{-1}\Fs_0$, is the same equation. (This is the hatted form of the identity $(z\,\dd/\dd v)^3v(z\,\dd/\dd v)^2=v(z\,\dd/\dd v)^5+3z(z\,\dd/\dd v)^4$, which follows from $[z\,\dd/\dd v,v]=z$: the leading part of $\PFfive$ at $u=0$ is the operator of the model.) By the uniqueness half of \cref{prop:set-modrows}, or by the pivot above, $\Phi_0=\Fs_0$. The free coefficient $\fc_m$ sits at $u^{4m}$ with $m\geq 1$, never at $u^0$, so the conclusion holds for every choice of the free coefficients.
\end{proof}

\begin{rmk}[The model does not fix the normalisation]
\label{rmk:con-model}
By \cref{lem:con-u0} the conifold value $\Phi(\zh,0)=\Fs_0(\zh)$ is the same for every choice of the free coefficients: the model alone does not fix the normalisation. The free coefficient $\fc_m$ sits at $u^{4m}$; by \cref{prop:con-deg} it is the top coefficient of $\mathrm{ph}_{5m}$, that is, the value of $L^{-1}A_{5m}$ at $L=0$. At the conifold point the normalisation is fixed by the gauge: $\Phican$ is singled out by the canonical gauge through (G) (\cref{thm:gau-G,def:gau-can}), not by the model.
\end{rmk}

% Section 8 (part 2 of 8): the degree bound and the three points.

\begin{prop}[The degree bound]
\label{prop:con-deg}
Let $(\fc_m)_{m\geq1}$ be arbitrary rational numbers. There is exactly one solution $\Phi=\sum_n\mathrm{ph}_n\zh^n$ of HPF with $\mathrm{ph}_0=1$, every $\mathrm{ph}_n$ in $\bbQ\fps{u}$ and $[u^{4m}]\mathrm{ph}_{5m}=\fc_m$ for every $m\geq1$. For it every $\mathrm{ph}_n$ is a polynomial of degree at most $\lfloor4n/5\rfloor$, and every resonance condition of \cref{lem:gau-HPF} holds trivially: $\deg r_n<4n/5$, so that $[u^{4m}]r_{5m}=0$. Equivalently, every $A_n$ lies in $L\bbQ[L]$ and has degree at most $4n+1$, and for $n\geq1$ the value of $L^{-1}A_n$ at $L=0$ is $\fc_{n/5}$ if $5\mid n$ and $0$ otherwise ($A_0=L$). Every solution of HPF with $\mathrm{ph}_0=1$ and coefficients in $\bbQ\fps{u}$ is one of these. The same holds in every colour.
\end{prop}

The proposition is not used by the proof of the main theorem. Its proof uses neither Theorem~U nor \cref{thm:con-Pa}, and it constructs a polynomial solution of HPF for every choice of the free coefficients. We do not claim that the bound $\lfloor4n/5\rfloor$ is attained.

\begin{proof}
Colour $i$ reduces to colour $0$: its exponent $u_i$ has $\Dq u_i=\zeta^iL$ (\cref{cor:set-cyclicPF5}), and $L':=\zeta^iL$ satisfies $\Dq L'=(L'^6-L')/5$ and $L'^{-5}=L^{-5}$, so the argument below applies verbatim in the variable $L'$. We work in colour $0$. The letters $\theta_L$, $\mathfrak p$, $\cP_j$, $\mathfrak e$, $R_n$, $\tilde r_n$, $\mathcal M_1$, $\mathcal M_{-4}$ are local to this proof. Let $\theta_L:=L\,\dd/\dd L$; on $\bbQ[L,L^{-1}]$ the derivation $\Dq$ is $-\frac15(1-L^5)\theta_L$, because $\Dq L=(L^6-L)/5$, and $q=(1-L^{-5})/5^5$. Put $\mathfrak p:=5^{-5}\prod_{k=1}^5\bigl(5(z\Dq+L)+kz\bigr)$; then
\begin{equation}
\label{eq:con-PL}
\PFfive^{L}=(z\Dq+L)^5-1+L^{-5}(1-L^5)\,\mathfrak p=\sum_{j=0}^{5}z^j\cP_j ,
\end{equation}
with operators $\cP_j$ on $\bbQ[L,L^{-1}]$, polynomials in $\theta_L$ with coefficients in $\bbQ[L,L^{-1}]$. For the flat amplitude $\psi:=\sum_nA_nz^n=s^{-1}\Phi$ the equation $[z^{n+1}]\PFfive^{L}\psi=0$ is the $\zh^{n+1}$-coefficient of HPF, by \eqref{eq:con-chain}.

\emph{Step 1 (the term $z^0$).} At $z^0$ one has $z\Dq+L=L$ and $\mathfrak p=L^5$, so $\cP_0=L^5-1+L^{-5}(1-L^5)L^5=0$.

\emph{Step 2 (the left factor).} $z\Dq=(1-L^5)\circ(-\tfrac z5\theta_L)$, and multiplications by functions of $L$ commute. A term of positive $z$-degree of $(z\Dq+L)^5$ is a word in $L$ and $z\Dq$ containing $z\Dq$; the factors to the left of its leftmost $z\Dq$ are multiplications by $L$, so the word has the left factor $1-L^5$. The term $L^{-5}(1-L^5)\mathfrak p$ has it explicitly. With Step 1: for $j\geq1$, $\cP_j=(1-L^5)\circ\tilde\cP_j$ for an operator $\tilde\cP_j$ on $\bbQ[L,L^{-1}]$.

\emph{Step 3 (valuation).} For $m\geq1$, $\bigl(5(z\Dq+L)+kz\bigr)L^m=5L^{m+1}+z(k-m)L^m+zmL^{m+5}$. The five factors of $\mathfrak p$ commute; apply first those with $k=m,m+1,\dots,5$ (if $m\leq5$): each kills the lowest power when that power equals $k$, and no factor lowers the valuation. Hence $\mathfrak p(L^m)\in L^6\bbQ[L][z]$ for $m\geq1$. Since $(z\Dq+L)L^m=L^{m+1}-\frac{zm}5(1-L^5)L^m$ does not lower the valuation either, every $\cP_j$ maps $L\bbQ[L]$ into itself.

\emph{Step 4 (the first-order part).} On $L^m$ the $z^1$-part of $(z\Dq+L)^5$ is $\sum_{i=0}^4L^i(z\Dq)L^{m+4-i}=-(m+2)z(1-L^5)L^{m+4}$, and that of $\mathfrak p$ is $z\bigl(3-(m+2)(1-L^5)\bigr)L^{m+4}$, the $3$ being $\sum_kk/5$. Hence $\cP_1L^m=(1-m)(1-L^5)L^{m-1}$, that is, $\cP_1=(1-L^5)\circ\mathfrak e$ with $\mathfrak e(L^m):=(1-m)L^{m-1}$.

\emph{Step 5 (weights).} Give $L$ the weight $1$, $z$ the weight $-4$ and $\theta_L$ the weight $0$. Then $z\Dq+L=\mathcal M_1+\mathcal M_{-4}$ with $\mathcal M_1:=L+\tfrac z5L^5\theta_L$ of weight $1$ and $\mathcal M_{-4}:=-\tfrac z5\theta_L$ of weight $-4$, and $5(z\Dq+L)+kz=5\mathcal M_1+(5\mathcal M_{-4}+kz)$, the second part of weight $-4$. A product of five factors one of which has weight $-4$ has weight at most $0$. Hence $(z\Dq+L)^5$ and $\mathfrak p$ have the same part $\mathcal M_1^5$ of weight $5$, and all their other parts have weight at most $0$. In \eqref{eq:con-PL} the parts of weight $5$ cancel, and $1$ and $L^{-5}\mathfrak p$ have weight at most $0$. So every term $z^jL^a\theta_L^b$ of $\PFfive^{L}$ has $a\leq4j$.

\emph{Step 6 (induction).} We show: for every $n$ there is exactly one $A_n$ in $L^{4n+1}\bbQ\fps{L^{-5}}$ with $[z^{n+1}]\PFfive^{L}\psi=0$ and $[L^1]A_{5m}=\fc_m$, and it lies in $L\bbQ[L]$ and has degree at most $4n+1$. For $n=0$, $A_0=L$. Assume this for $A_0,\dots,A_{n-1}$. By Step 1 the equation is $\cP_1A_n=R_n:=-\sum_{j=2}^5\cP_jA_{n+1-j}$. By Step 3, $R_n\in L\bbQ[L]$. By Step 5, applied to the terms $z^kA_k$ ($k<n$), which have weight at most $1$, $z^{n+1}R_n$ has weight at most $1$, so $\deg R_n\leq4n+5$. By Step 2, $R_n\in(1-L^5)\bbQ[L,L^{-1}]$. As $1-L^5$ is a unit of $\bbQ\fps{L}$, $\tilde r_n:=R_n/(1-L^5)$ lies in $L\bbQ\fps{L}\cap\bbQ[L,L^{-1}]=L\bbQ[L]$, and $\deg\tilde r_n\leq4n$. By Step 4, and since $1-L^5$ is not a zero divisor in $\bbQ\lser{L^{-1}}$, the equation reads $\mathfrak eA_n=\tilde r_n$, that is,
\[
(1-m)\,[L^m]A_n=[L^{m-1}]\tilde r_n\qquad(m\in\bbZ).
\]
At $m=1$ this says $0=[L^0]\tilde r_n$, which holds because $\tilde r_n\in L\bbQ[L]$: the resonance is trivial. For $m\neq1$ the equation determines $[L^m]A_n$, which vanishes for $m\leq0$ and for $m>4n+1$. Finally $\PFfive^{L}$ maps $L^mz^j$ to combinations of monomials $L^{m'}z^{j'}$ with $m'+j'\equiv m+j\bmod5$, and $\psi$ has only monomials with $m+j\equiv1$; so $A_n$ has only exponents $\equiv4n+1\bmod5$, and the coefficient $[L^1]A_n$, which the equation leaves free, occurs only for $5\mid n$, where it is $[u^{4n/5}]\mathrm{ph}_n$. This proves the claim.

\emph{Step 7 (translation).} $A_n=\sum_k[u^k]\mathrm{ph}_n\,L^{4n+1-5k}$ lies in $L\bbQ[L]$ if and only if $[u^k]\mathrm{ph}_n=0$ whenever $4n+1-5k\leq0$, that is, if and only if $\deg\mathrm{ph}_n\leq\lfloor4n/5\rfloor$; and $L^{-1}A_n$ at $L=0$ is $[L^1]A_n$. Comparing $[z^{n+1}]\PFfive^{L}\psi=(1-L^5)(\mathfrak eA_n-\tilde r_n)$ with the $\zh^{n+1}$-coefficient $(1-u)[(4n-5\theta_u)\mathrm{ph}_n+r_n]$ of HPF through \eqref{eq:con-chain} gives $\tilde r_n=L^{4n}r_n(L^{-5})$; so $\tilde r_n\in L\bbQ[L]$ means $\deg r_n<4n/5$, and $[L^0]\tilde r_{5m}=[u^{4m}]r_{5m}$ is the resonance coefficient of \cref{lem:gau-HPF}. The $u$-recursion of \cref{lem:gau-HPF} determines $\mathrm{ph}_n$ from $\mathrm{ph}_0,\dots,\mathrm{ph}_{n-1}$ and $\fc_{n/5}$; hence every solution with $\mathrm{ph}_0=1$ and coefficients in $\bbQ\fps{u}$ is the one constructed for its free coefficients.
\end{proof}

\begin{rmk}[Three points]
\label{rmk:con-points}
By \cref{prop:con-deg}, for a colour-$0$ WKB function solving HPF with coefficients in $\bbQ\fps{u}$ the series $L^{-1}\sum_nA_nz^n$ lies in $\bbQ[L]\fps{z}$. Such series connect three points. At $L=1$ ($q=0$, $u=1$, $\zh=z$) the value of the actual amplitude, whose coefficients are Laurent polynomials in $L$ (\cref{thm:gau-R0}~(iii), \cref{thm:con-glue}), is prescribed by (QRR): it is $h(z)$ (\cref{thm:set-qrr}). At $L=0$ ($q=\infty$) the value is $1+\sum_m\fc_mz^{5m}$; for the full-Stirling series of \cref{sec:con-form} it is determined in \cref{thm:gam-ginf}. At $L=\infty$ ($u=0$, the conifold point) the leading coefficients are those of the model (\cref{lem:con-u0}), and the normalisation is fixed by the canonical gauge, through the top coefficients $\fc_m$ (\cref{thm:con-Pb,thm:con-glue}).
\end{rmk}

% Section 8 (part 3 of 8): the full-Stirling series, its rationality, the 1-form and PhiSP.

\subsection{The full-Stirling series, its rationality, the \texorpdfstring{$1$}{1}-form and \texorpdfstring{$\PhiSP$}{the full-Stirling WKB function}}
\label{sec:con-form}

\begin{defn}[The full-Stirling saddle series]
\label{def:gam-c}
Let $\Cstar:=\bbC\setminus\{0,1,\zeta,\zeta^2,\zeta^3,\zeta^4\}$ and $q(L):=(1-L^{-5})/5^5$; no fifth root is taken. For $L\in\Cstar$ fix an open disc $\bbD_L$ around $L$ containing neither $0$ nor any $\zeta^j$, holomorphic branches of $\log(5\sigma)$ and of $\log(\zeta^j-\sigma)$, $j=0,\dots,4$, on $\bbD_L$, and put
\[
\lambda:=-5\log(5L)+\sum_{j=0}^4\log(\zeta^j-L),\qquad\text{so that}\qquad \re^{\lambda}=\frac{\prod_j(\zeta^j-L)}{(5L)^5}=-q(L)
\]
(the letter $\lambda$ is local to this section and \cref{sec:gam}). On $\bbD_L$ define
\begin{align*}
\mathsf S_0(\sigma)&:=\lambda\sigma+5\sigma\log(5\sigma)-5\sigma+\sum_{j=0}^4\bigl[(\zeta^j-\sigma)\log(\zeta^j-\sigma)-(\zeta^j-\sigma)\bigr],\\
\mathsf S_1(\sigma)&:=\tfrac12\log(5\sigma)-\tfrac12\sum_{j=0}^4\log(\zeta^j-\sigma),\qquad
\mathsf S_{2k}(\sigma):=\bk_k\Bigl[(5\sigma)^{1-2k}+\sum_{j=0}^4(\zeta^j-\sigma)^{1-2k}\Bigr]\quad(k\geq1),\\
A_{\mathrm{St}}(\sigma,z)&:=\exp\Bigl(\mathsf S_1(\sigma)+\sum_{k\geq1}z^{2k-1}\mathsf S_{2k}(\sigma)\Bigr).
\end{align*}
Then $\mathsf S_0'(L)=0$ and $\mathsf S_0''=5/\sigma+\sum_j(\zeta^j-\sigma)^{-1}=5/(\sigma(1-\sigma^5))$, which does not vanish at $L$. The \emph{full-Stirling saddle series} is the normalised stationary-phase series (\cref{def:fsp-wick}), with constant term $1$, in the variable $\sigma$ with the measure $\dd\sigma$:
\[
c(z;L):=\SPh_L[\mathsf S_0;A_{\mathrm{St}}](z)=\sum_{n\geq0}c_n(L)z^n,\qquad c_0(L)=1 .
\]
For colour $i$ the series is defined in the same way at the critical point $\sigma=\zeta^iL$, with the logarithm of $-q$ chosen so that $\zeta^iL$ is critical.
\end{defn}

With $\log(a/z):=\log a-\log z$ and $\Gst(1+w):=w\,\Gst(w)$ (\cref{def:fsp-stirling}) one has, identically in $\sigma\in\bbD_L$,
\[
\re^{\lambda\sigma/z}\,\Gst(1+5\sigma/z)\prod_{j=0}^4\Gst\bigl((\zeta^j-\sigma)/z\bigr)=(2\pi)^3z^2\exp\Bigl(\mathsf S_0/z+\mathsf S_1+\sum_kz^{2k-1}\mathsf S_{2k}\Bigr),
\]
because the coefficient of $(\log z)/z$ is $-[5\sigma+\sum_j(\zeta^j-\sigma)]=-\sum_j\zeta^j=0$ and that of $\log z$ is $-\frac12+\frac52=2$. So $c(z;L)$ is the normalised saddle series at $\sigma=L$ of the Mellin--Barnes summand $\Gamma(1+5\sigma/z)\prod_j\Gamma((\zeta^j-\sigma)/z)(-q)^{\sigma/z}\,\dd\sigma$, with all six Gamma factors replaced by $\Gst$ and $(-q)^{\sigma/z}:=\re^{\lambda\sigma/z}$. In the Euler integrals, $\Gamma(1+5\sigma/z)=\int y^{5\sigma/z}\re^{-y}\,\dd y$ carries the measure $\dd y$ and $\Gamma(w)=\int x^w\re^{-x}\,\dd x/x$ the measure $\dd x/x$; this matters only in \cref{sec:gam}.

\emph{Taylor data.} Consider the local power sums $p_k(L):=\sum_j(\zeta^j-L)^{-k}$. Their generating function is
\[
\sum_{k\geq1}p_k(L)\,t^{k-1}=\sum_j(\zeta^j-L-t)^{-1}=\frac{5(L+t)^4}{1-(L+t)^5},
\]
so that $p_k\in(1-L^5)^{-k}\bbQ[L]$, and one has
\begin{gather*}
\mathsf S_0(L+t)-\mathsf S_0(L)=\sum_{k\geq1}\frac{5(-1)^{k+1}L^{-k}+p_k(L)}{k(k+1)}\,t^{k+1},\qquad
-1/\mathsf S_0''(L)=-L(1-L^5)/5,\\
\mathsf S_1(L+t)-\mathsf S_1(L)=\sum_{k\geq1}\frac{(-1)^{k+1}L^{-k}+p_k(L)}{2k}\,t^{k},\\
[t^r]\,\mathsf S_{2k}(L+t)=\bk_k\binom{1-2k}{r}\bigl[5^{1-2k}L^{1-2k-r}+(-1)^rp_{2k-1+r}(L)\bigr].
\end{gather*}

\begin{prop}[Rationality]
\label{prop:gam-rational}
The series $c(z;L)$ does not depend on the disc $\bbD_L$ nor on the branches of the logarithms, and every $c_n$ is the value at $L$ of an element of $\AL=\bbQ[L,L^{-1},(L^5-1)^{-1}]$. More precisely, $c(z;L)$ depends only on the Taylor coefficients of order at least $2$ of the phase at $L$ and on the normalised amplitude $A_{\mathrm{St}}/A_{\mathrm{St}}(L,0)$; $\lambda$ and the sign convention for $q$ never enter it, they only fix the dictionary $L\leftrightarrow q$.
\end{prop}

\begin{proof}
(1) $\mathsf S_0$ and $A_{\mathrm{St}}$ are holomorphic on $\bbD_L$, $\mathsf S_0'(L)=0$ and $\mathsf S_0''(L)=5/(L(1-L^5))\neq0$. (2) By \cref{lem:fsp-univ}, $c_n(L)$ is a polynomial over $\bbQ$ in the Taylor coefficients at $L$ of $\mathsf S_0$ of orders $3,\dots,2n+2$, in those of $\mathsf S_1-\mathsf S_1(L)$ and of the $\mathsf S_{2k}$ (which determine the normalised amplitude), and in $-1/\mathsf S_0''(L)$. (3) These are the expressions displayed above: polynomials in $L^{-1}$, in the $p_k\in(1-L^5)^{-k}\bbQ[L]$ and in $-L(1-L^5)/5$. (4) Hence $c_n\in\AL$. (5) The Taylor coefficients of order at least $1$ of $\mathsf S_0$, $\mathsf S_1$, $\mathsf S_{2k}$ are rational in $\sigma$ and carry no branch, and $\exp(\mathsf S_1(\sigma)-\mathsf S_1(L))$ is the branch equal to $1$ at $L$. Replacing $\log(5\sigma)$ by $\log(5\sigma)+2\pi\ri n$ and $\log(\zeta^j-\sigma)$ by $\log(\zeta^j-\sigma)+2\pi\ri n_j$ changes $\mathsf S_0'$ by $2\pi\ri(5n-\sum_jn_j)$, which the formula for $\lambda$ compensates; then $\mathsf S_0$ changes by the constant $2\pi\ri\sum_j\zeta^jn_j$ and $\mathsf S_1$ by a constant, which the normalised series does not see (\cref{lem:fsp-const}). A logarithm $\lambda+2\pi\ri k$ of $-q$ with $k\neq0$ gives $\mathsf S_0'(L)=2\pi\ri k\neq0$, so $L$ is critical only for $\lambda$; a smaller disc changes nothing. So $c(z;L)$ depends only on $L$, and only through the data named in the statement.
\end{proof}

\begin{lem}[One $1$-form]
\label{lem:con-form}
Let $\Xi:=-\log(1-u)\in u\bbQ\fps{u}$ and $\cG(x):=\sum_{m\geq1}x^{5m-1}/(m(5m-1))$.
\begin{enumerate}
\item[(a)] In the colour-$0$ summand
\[
q^{\sigma/z}\,\Gamma(1+5\sigma/z)\Big/\prod_{j=0}^4\Gamma\bigl(1+(\sigma-\zeta^j)/z\bigr)
\]
expand $\Gamma(1+5\sigma/z)$ and the five factors $1/\Gamma(1+(\sigma-\zeta^j)/z)$ by Stirling's series, and put $x=1/\sigma$. The exponent is $S(\sigma)/z$ with $S(\sigma):=\sigma\log q+5\sigma\log(5\sigma)-\sum_j(\sigma-\zeta^j)\log(\sigma-\zeta^j)$: the terms in $(\log z)/z$ and the terms linear in $\sigma/z$ cancel. Moreover $S(1/x)=-\cW(x)$ with $\cW(x):=\Xi/x+\cG(x)$, and $\dd\cW/\dd x=-(\Xi+\log(1-x^5))/x^2$; the critical points are $x^5=u$.
\item[(b)] The half-densities and the measure combine to $(5\sigma/(\sigma^5-1))^{1/2}\,\dd\sigma=-\sqrt5\,(1-x^5)^{-1/2}\,\dd x$. Up to a factor independent of $x$, the summand times $\dd\sigma$ is $\re^{-\cW/z}\fa\,\dd x$ with
\[
\log\fa(x,z)=-\tfrac12\log(1-x^5)+\sum_{k\geq1}\bk_k(xz)^{2k-1}\Bigl[5^{1-2k}-5\sum_{P\geq0}\binom{2k-2+5P}{5P}x^{5P}\Bigr];
\]
$\fa$ does not depend on $u$, $\fa(0,z)=1$, and the coefficient of $x^a$ in $\fa$ is a polynomial in $z$ of degree at most $a$.
\item[(c)] In the hatted variables $x=sX$, $z=s^4\zh$ one has $\cW/z=\hat\cW/\zh$ with
\[
\hat\cW(X;u)=\sum_{m\geq1}\frac{u^{m-1}}m\Bigl[\frac1X+\frac{X^{5m-1}}{5m-1}\Bigr],\qquad \hat\cW'(X)=-\frac1{uX^2}\log\frac{1-uX^5}{1-u};
\]
$X=1$ is critical for every $u$, $\hat\cW''(1)=5/(1-u)$, $\hat\cW(1)=\frac54\tau(u)^{-1}$, and $[t^k]\hat\cW(1+t)=\sum_{m\geq1}u^{m-1}[(-1)^k+\binom{5m-1}{k}/(5m-1)]/m$ lies in $\bbQ\fps{u}$. The hatted amplitude $\hat\fa(X;\zh):=\fa(sX,s^4\zh)$ has
\[
\log\hat\fa=-\tfrac12\log(1-uX^5)+\sum_{k\geq1}\bk_k(uX\zh)^{2k-1}\Bigl[5^{1-2k}-5\sum_{P\geq0}\binom{2k-2+5P}{5P}u^PX^{5P}\Bigr];
\]
it lies in $\bbQ[X,\zh]\fps{u}$, and $\hat\fa(1;0)=(1-u)^{-1/2}$.
\item[(d)] The series $c(z;L)$ of \cref{def:gam-c} (the first Mellin--Barnes summand) and the normalised saddle series at $\sigma=L$ of the second summand of (a), with all Gamma factors replaced by their Stirling forms and the logarithm of $q$ chosen so that $L$ is critical, coincide.
\end{enumerate}
\end{lem}

% Section 8 (part 4 of 8): proof of the 1-form lemma; PhiSP and the residue; Theorem P(a), statement and part (i).

\begin{proof}
(a) Replace each $\Gamma(1+y)$ by its Stirling form (\cref{def:fsp-stirling})
\[
\Gst(1+y)=y\,\Gst(y)=\sqrt{2\pi}\,y^{y+1/2}\re^{-y}\exp\Bigl(\sum_{k\geq1}\bk_ky^{1-2k}\Bigr),
\]
with $y=5\sigma/z$ for the numerator and $y=(\sigma-\zeta^j)/z$ for the factors of the denominator, and $\log(a/z)=\log a-\log z$. The terms in $(\log z)/z$ are $-(\log z/z)\,[5\sigma-\sum_j(\sigma-\zeta^j)]=0$, those linear in $\sigma/z$ are $[-5\sigma+\sum_j(\sigma-\zeta^j)]/z=0$, and $q^{\sigma/z}=\re^{\sigma\log q/z}$; what remains of order $1/z$ is $S(\sigma)/z$. At $\sigma=1/x$ write $\log(1/x-\zeta^j)=\log(1-\zeta^jx)-\log x$; with $\sum_j(1/x-\zeta^j)=5/x$ and $\log q+5\log5=\log(1-u)$ this gives $S(1/x)=-\Xi/x-\sum_j(1/x-\zeta^j)\log(1-\zeta^jx)$. Expanding $\log(1-\zeta^jx)=-\sum_n\zeta^{jn}x^n/n$ and summing over $j$, only the terms with $5\mid n$ survive in $\sum_j\log(1-\zeta^jx)$ and only those with $5\mid n+1$ in $\sum_j\zeta^j\log(1-\zeta^jx)$, and
\[
\sum_j(1/x-\zeta^j)\log(1-\zeta^jx)=-\sum_{m\geq1}\frac{x^{5m-1}}m+5\sum_{m\geq1}\frac{x^{5m-1}}{5m-1}=\cG(x).
\]
Then $\cG'(x)=\sum_mx^{5m-2}/m=-x^{-2}\log(1-x^5)$ gives $\dd\cW/\dd x$, which vanishes if and only if $\log(1-x^5)=\log(1-u)$.

(b) The half-densities are $(5\sigma/z)^{1/2}$ from $\Gamma(1+5\sigma/z)$ and $((\sigma-\zeta^j)/z)^{-1/2}$ from the denominator; their product is $z^2(5\sigma/\prod_j(\sigma-\zeta^j))^{1/2}=z^2(5\sigma/(\sigma^5-1))^{1/2}$, and at $\sigma=1/x$ one has $5\sigma/(\sigma^5-1)=5x^4/(1-x^5)$ and $\dd\sigma=-\dd x/x^2$. The Stirling series contribute $\sum_k\bk_k[(5\sigma/z)^{1-2k}-\sum_j((\sigma-\zeta^j)/z)^{1-2k}]=\sum_k\bk_k(xz)^{2k-1}[5^{1-2k}-\sum_j(1-\zeta^jx)^{1-2k}]$, and $\sum_j(1-\zeta^jx)^{-r}=5\sum_{5\mid i}\binom{r-1+i}{i}x^i$ (expand and sum over $j$). The constant $-(2\pi)^{-2}\sqrt5$ and the power $z^2$ do not depend on $x$; what is left is $\re^{-\cW/z}\fa\,\dd x$. Every monomial of $\log\fa$ has $z$-degree at most its $x$-degree, hence so has every monomial of $\fa$.

(c) With $x=sX$ and $z=s^4\zh$ one has $\Xi/(xz)=(\Xi/u)/(X\zh)$ with $\Xi/u=\sum_mu^{m-1}/m$, and $\cG(sX)/(s^4\zh)=\sum_mu^{m-1}X^{5m-1}/(m(5m-1)\zh)$; this is $\hat\cW/\zh$. Termwise in $u$,
\[
\hat\cW'(X)=\sum_m\frac{u^{m-1}}m(X^{5m-2}-X^{-2})=(uX^2)^{-1}\sum_m\frac{u^m}m(X^{5m}-1),
\]
which is the formula stated and vanishes at $X=1$; differentiating the logarithm at $X=1$ gives $\hat\cW''(1)=5/(1-u)$; and $\hat\cW(1)=5\sum_mu^{m-1}/(5m-1)=\frac54\sum_{n\geq0}4u^n/(5n+4)=\frac54\tau^{-1}$ (\cref{def:set-tau}). The Taylor coefficients at $X=1$ follow from those of $(1+t)^{-1}$ and $(1+t)^{5m-1}$. Finally $xz=uX\zh$ and $x^5=uX^5$ in (b).

(d) By \cref{lem:fsp-univ,lem:fsp-const}, a normalised saddle series depends only on the Taylor coefficients of order at least $2$ of the phase at the critical point and on the normalised amplitude. For $y_j:=(\zeta^j-\sigma)/z$ the formal reflection identity \cref{lem:fsp-refl} gives $\log\Gst(y_j)+\log\Gst(1-y_j)=\log2\pi+(y_j-\frac12)(\log y_j-\log(-y_j))$, and $\log y_j-\log(-y_j)=\pi\ri\,\epsilon_j$ with an odd integer $\epsilon_j$ fixed by the branches. Since $1-y_j=1+(\sigma-\zeta^j)/z$, the product $\prod_j\Gst(1+(\sigma-\zeta^j)/z)^{-1}$ is a constant times $\re^{c/z}\re^{\pi\ri(\sum_j\epsilon_j)\sigma/z}\prod_j\Gst((\zeta^j-\sigma)/z)$, with $c$ independent of $\sigma$. Moreover $q^{\sigma/z}=\re^{(\lambda+\pi\ri\epsilon)\sigma/z}$ with an odd integer $\epsilon$, because the logarithms of $q$ and of $-q$ differ by an odd multiple of $\pi\ri$. So the second summand, in Stirling form, is a constant times $\re^{c/z}$ times the first with $\lambda$ replaced by $\lambda+\pi\ri e$, where $e:=\epsilon+\sum_j\epsilon_j$ is even. If $L$ is critical for both, then $e=0$ (step (5) of the proof of \cref{prop:gam-rational}). The two integrands then have the same phase up to an additive constant and the same amplitude up to a constant factor, so their normalised series coincide. (Directly: the phase $S$ of (a) has $S''=5/\sigma-\sum_j(\sigma-\zeta^j)^{-1}=\mathsf S_0''$, so $S-\mathsf S_0$ is constant near $L$ once both are critical at $L$; the half-densities $(5\sigma)^{1/2}\prod_j(\sigma-\zeta^j)^{-1/2}$ and $\re^{\mathsf S_1}$ differ by a constant factor; and the Stirling parts agree, because $(\zeta^j-\sigma)^{1-2k}=-(\sigma-\zeta^j)^{1-2k}$.)
\end{proof}

\begin{defn}[The full-Stirling WKB function and the residue]
\label{def:con-PhiSP}
The \emph{full-Stirling WKB function} is the normalised stationary-phase series, with constant term $1$, of $\re^{-\hat\cW/\zh}\hat\fa\,\dd X$ at $X=1$:
\[
\PhiSP(\zh,u):=\SPh_{X=1}[-\hat\cW;\hat\fa](\zh)=\sum_{n\geq0}\mathrm{ph}^{\mathrm{SP}}_n(u)\zh^n,\qquad\mathrm{ph}^{\mathrm{SP}}_0=1 .
\]
Equivalently,
\[
\PhiSP=\mathrm{SP}_{X=1}\bigl[\re^{-\hat\cW/\zh}\hat\fa\,\dd X\bigr]\Big/\bigl((2\pi\zh/5)^{1/2}\re^{-\hat\cW(1)/\zh}\bigr),
\]
the full stationary phase of \cref{def:fsp-wick} divided by $(2\pi\zh/5)^{1/2}\re^{-\hat\cW(1)/\zh}$, which is the Gaussian factor $(2\pi\zh)^{1/2}\hat\cW''(1)^{-1/2}$ times the critical exponential $\re^{-\hat\cW(1)/\zh}$ times the amplitude value $\hat\fa(1;0)=(1-u)^{-1/2}$; in every comparison below this $\zh$-dependent factor is common to both sides. By \cref{lem:con-form}~(c) and \cref{lem:fsp-univ}, every $\mathrm{ph}^{\mathrm{SP}}_n$ lies in $\bbQ\fps{u}$. The \emph{residue} of the flat $1$-form is
\[
\cV:=\Res_{x=0}\re^{-\cW/z}\fa\,\dd x:=\sum_{k\geq1}\frac{(-\Xi/z)^k}{k!}\,[x^{k-1}]\bigl(\re^{-\cG/z}\fa\bigr).
\]
Since $\Xi\in u\bbQ\fps{u}$, $\cG\in x^4\bbQ\fps{x}$ and the coefficient of $x^a$ in $\fa$ has $z$-degree at most $a$, each $u$-order of $\cV$ is a finite sum, and $\cV\in z^{-1}\bbQ[z^{-1}]\fps{u}$.
\end{defn}

\subsection{Theorem P}
\label{sec:con-P}

The ring $\AL$ carries the derivation $\Dq$ with $\Dq L=(L^6-L)/5$; it contains $q=q(L)$, and $\Dq q=q$, so that $\Dq=q\,\dd/\dd q$ on $\AL$. Let $\jmath\colon\AL\to\bbQ\lser{s}$ be the ring homomorphism with $\jmath(L)=s^{-1}$; it is injective, because $\bbQ(L)\to\bbQ\lser{s}$, $L\mapsto s^{-1}$, is a field embedding.

\begin{thm}[Theorem P(a)]
\label{thm:con-Pa}
Let $c(z;L)=\sum_nc_n(L)z^n$ be the full-Stirling saddle series, with $c_n\in\AL$ (\cref{def:gam-c,prop:gam-rational}).
\begin{enumerate}
\item $\PFfive^{L}\bigl(L\,c(z;L)\bigr)=0$ in $\AL\fps{z}$.
\item In $\bbQ\lser{s}\fps{z}$ one has $\jmath(c)(z)=\PhiSP(s^{-4}z;s^5)$; that is, $\jmath(c_n)=s^{-4n}\mathrm{ph}^{\mathrm{SP}}_n$ for every $n$.
\item $\PhiSP$ solves HPF in $\bbQ\fps{u}\fps{\zh}$.
\item Every $\mathrm{ph}^{\mathrm{SP}}_n$ is a polynomial in $u$; hence $c_n(L)=L^{4n}\mathrm{ph}^{\mathrm{SP}}_n(L^{-5})$ in $\AL$, and $c_n\in\bbQ[L,L^{-1}]$.
\end{enumerate}
\end{thm}

The proof uses \cref{def:gam-c,prop:gam-rational,lem:con-form,lem:gau-HPF,lem:gau-poly} and, from \cref{sec:fsp}, \cref{lem:fsp-univ,lem:fsp-const,lem:fsp-ift,lem:fsp-generic,lem:fsp-dq,lem:fsp-transl,lem:fsp-stirfe,lem:fsp-cov,lem:fsp-homog}. It uses neither Theorem~U, nor the cyclic form \cref{thm:gau-cyc}, nor (QRR), nor \cref{prop:con-deg}, nor Theorem Gamma; part (ii) does not use (iv).

\begin{proof}[Proof of (i)]
In this proof $\mathsf S_0$ is written as a formal phase at $\sigma=L$, that is, by its Taylor data at $L$, which lie in $\AL$ (\cref{def:gam-c}), and $\sigma=L+t$. For $\chi\in\bbQ[\sigma,z]$ put $c\langle\chi\rangle(z;L):=\SPn_L[\mathsf S_0;\chi A_{\mathrm{St}}]/A_{\mathrm{St}}(L,0)$, so that $c=c\langle1\rangle$; by \cref{def:fsp-wick} the map $\chi\mapsto c\langle\chi\rangle$ extends $\AL[z]$-linearly to $\AL[\sigma,z]$. The letters $\epsilon$, $F_\epsilon$, $\sigma_\epsilon$, $\chi$, $a$, $b_z$ are local to this proof.

\emph{Step 1 ($\Dq$ on saddle series).} Let $\epsilon$ be a formal parameter and $F_\epsilon(\sigma):=\mathsf S_0(\sigma)+\epsilon\sigma$. By \cref{lem:fsp-generic}, which rests on \cref{lem:fsp-ift}, $F_\epsilon$ has a unique critical point $\sigma_\epsilon\in L+\epsilon\AL\fps{\epsilon}$, and the ring homomorphism $\AL\to\AL\fps{\epsilon}$, $L\mapsto\sigma_\epsilon$, is $\exp(\epsilon\Dq)=\sum_m\epsilon^m\Dq^m/m!$ (both $\sigma_\epsilon$ and $\exp(\epsilon\Dq)(L)$ solve $y'=y(y^5-1)/5$, $y(0)=L$, since $\sigma_\epsilon'=-1/\mathsf S_0''(\sigma_\epsilon)$). The Taylor coefficients at $\sigma_\epsilon$ of $F_\epsilon$ of order at least $2$, and those of $\chi A_{\mathrm{St}}/A_{\mathrm{St}}(\sigma_\epsilon,0)$ for $\chi\in\bbQ[\sigma,z]$, are values at $\sigma_\epsilon$ of elements of $\bbQ[\sigma,\sigma^{-1},(\sigma^5-1)^{-1}][z]$ (the derivatives of $\mathsf S_0$, $\mathsf S_1$ and the $\mathsf S_{2k}$ are rational), hence the images under $\exp(\epsilon\Dq)$ of the corresponding coefficients at $L$; this is where the coefficients of $\chi$ have to be $\Dq$-constants. By \cref{lem:fsp-univ},
\[
\SPn_{\sigma_\epsilon}[F_\epsilon;\chi A_{\mathrm{St}}]/A_{\mathrm{St}}(\sigma_\epsilon,0)=\exp(\epsilon\Dq)\bigl(c\langle\chi\rangle\bigr).
\]

\emph{Step 2 (the identity).} Apply \cref{lem:fsp-dq} over $\AL$, with the parameter $\epsilon$ at $\epsilon=0$, to the amplitude $\chi A_{\mathrm{St}}/A_{\mathrm{St}}(L,0)$, which does not depend on $\epsilon$; since $\partial_\epsilon F_\epsilon=\sigma$,
\[
\frac{\dd}{\dd\epsilon}\Big|_{\epsilon=0}\Bigl[\re^{(F_\epsilon(\sigma_\epsilon)-\mathsf S_0(L))/z}\Bigl(\frac{\mathsf S_0''(\sigma_\epsilon)}{\mathsf S_0''(L)}\Bigr)^{-1/2}\SPn_{\sigma_\epsilon}\bigl[F_\epsilon;\chi A_{\mathrm{St}}/A_{\mathrm{St}}(L,0)\bigr]\Bigr]=\frac1z\,c\langle\sigma\chi\rangle .
\]
By Step 1 the last factor in the bracket is $(A_{\mathrm{St}}(\sigma_\epsilon,0)/A_{\mathrm{St}}(L,0))\exp(\epsilon\Dq)(c\langle\chi\rangle)$. Now $A_{\mathrm{St}}(\sigma,0)^2=5\sigma/(1-\sigma^5)$ and $\mathsf S_0''(\sigma)=5/(\sigma(1-\sigma^5))$ have the quotient $\sigma^2$, and a series in $\epsilon$ with constant term $1$ has exactly one square root with constant term $1$; hence $(\mathsf S_0''(\sigma_\epsilon)/\mathsf S_0''(L))^{-1/2}A_{\mathrm{St}}(\sigma_\epsilon,0)/A_{\mathrm{St}}(L,0)=\sigma_\epsilon/L$, and the bracket is $\re^{(F_\epsilon(\sigma_\epsilon)-\mathsf S_0(L))/z}\exp(\epsilon\Dq)(L\,c\langle\chi\rangle)/L$. Since $\frac{\dd}{\dd\epsilon}F_\epsilon(\sigma_\epsilon)=\sigma_\epsilon$ ($\sigma_\epsilon$ is critical) and $\frac{\dd}{\dd\epsilon}|_{0}\exp(\epsilon\Dq)=\Dq$, the left side is $(zL)^{-1}(z\Dq+L)(L\,c\langle\chi\rangle)$. Hence
\begin{equation}
\label{eq:con-shiftD}
(z\Dq+L)\bigl(L\,c\langle\chi\rangle\bigr)=L\,c\langle\sigma\chi\rangle\quad\text{in }\AL\lser{z},\qquad\chi\in\bbQ[\sigma,z].
\end{equation}
The restriction on $\chi$ is necessary: for $\chi=q$ the two sides differ by $zqLc$, since $\Dq q=q$.

\emph{Step 3 ($\PFfive$).} $\PFfive=(z\Dq)^5-1-q\prod_{k=1}^5(5z\Dq+kz)$, with $q$ acting by multiplication on the left. Iterating \eqref{eq:con-shiftD} (the amplitudes $\sigma^m\chi$ stay in $\bbQ[\sigma,z]$) and using the $\AL[z]$-linearity of $\chi\mapsto c\langle\chi\rangle$,
\[
\PFfive^{L}(L\,c)=L\Bigl(c\langle\sigma^5-1\rangle-c\bigl\langle q\textstyle\prod_{k=1}^5(5\sigma+kz)\bigr\rangle\Bigr).
\]

\emph{Step 4 (the shift).} By \cref{lem:con-form}~(d) and its proof, the phase $S$ of the second Mellin--Barnes summand has the same Taylor coefficients at $L$ as $\mathsf S_0$ in every order at least $1$, and its amplitude is a constant multiple of $A_{\mathrm{St}}$; so $S$ may replace $\mathsf S_0$ in $c\langle\cdot\rangle$. Put $a:=(\sigma^5-1)A_{\mathrm{St}}/A_{\mathrm{St}}(L,0)$ and let $b_z(t):=\exp\bigl((S(L+t+z)-S(L+t))/z\bigr)\,a(L+t+z)$ be its translate by $z$ (\cref{lem:fsp-transl}~(b)); the exponent lies in $\AL\fps{t,z}$ without constant term, because $S'(L)=0$. We claim
\begin{equation}
\label{eq:con-bz}
b_z=q\prod_{k=1}^5(5\sigma+kz)\,A_{\mathrm{St}}/A_{\mathrm{St}}(L,0)\qquad\text{in }\AL\fps{t}\fps{z},\quad\sigma=L+t .
\end{equation}
Both sides lie in $\AL\fps{t}\fps{z}$ (the Taylor data of $A_{\mathrm{St}}/A_{\mathrm{St}}(L,0)$ at $L$ do, \cref{def:gam-c}), which embeds in $\AL(\zeta)\fps{t}\fps{z}$, $\AL(\zeta):=\AL\otimes\bbQ(\zeta)$, where $(L-\zeta^j)^{-1}$ exists; we prove \eqref{eq:con-bz} there, factor by factor, the product over $j$ being invariant under the Galois group of $\bbQ(\zeta)/\bbQ$. Write the Stirling-expanded second summand as the product of the six factors $\re^{\sigma\log q/z}$, $\Gst(1+5\sigma/z)$ and $\Gst(1+(\sigma-\zeta^j)/z)^{-1}$, $j=0,\dots,4$. Under $\sigma\mapsto\sigma+z$ the formal functional equation $\Gst(Y+1)=Y\,\Gst(Y)$ of \cref{lem:fsp-stirfe}, applied five times to the first Gamma factor and once to each of the others, with $1/Y=z/(5\sigma+kz)$ ($k=1,\dots,5$) and $1/Y=z/(\sigma+z-\zeta^j)$, all in $z\AL(\zeta)\fps{t}\fps{z}$, multiplies these factors by $q$, by $\prod_{k=1}^5(5\sigma/z+k)$ and by $z/(\sigma+z-\zeta^j)$ respectively; the terms in $\log z$ cancel before and after the shift, as in \cref{def:gam-c}. The product of the multipliers is $q\prod_k(5\sigma+kz)/((\sigma+z)^5-1)$, which is \eqref{eq:con-bz}. By \cref{lem:fsp-transl}~(b), $\SPn_L[S;b_z]=\SPn_L[S;a]$, that is, $c\langle q\prod_k(5\sigma+kz)\rangle=c\langle\sigma^5-1\rangle$; by Step 3, $\PFfive^{L}(L\,c)=0$.
\end{proof}

\input{sections/s08_connection_5}



